\section{Experiments}

%baseline : key baseline without modifficaiton
%add dataset + describe

\subsection{Experimental Design}

Our evaluation includes two primary CV tasks: \textbf{image classification} and \textbf{object detection}. We focus on measuring the trade-off between model accuracy and computational efficiency (both training and inference) while conducting ablation studies to isolate the impact of each architectural component.

\subsection{Task Specifications}

\paragraph{Image Classification}
We use the \textbf{ImageNet-100} dataset, a standard benchmark containing over 100 thousand images across 200 categories. Performance is measured using \textbf{Top-1 Accuracy}, which denotes the percentage of instances where the model's highest-probability prediction matches the ground-truth label.

\paragraph{Object Detection}
For localization tasks, we use \textbf{COCO} 2017 which is widely recognized for its complexity and variety of object scales. We report \textbf{Average Precision (AP)}, a metric that accounts for both the precision and recall of the predicted bounding boxes at various Intersection over Union (IoU) thresholds.

\subsection{Implementation Environment}
The framework is implemented in \textbf{Python} using \textbf{PyTorch}. All experiments, including training and inference profiling, are conducted on a local computer with a single NVIDIA RTX 5060Ti 16GB.

\subsection{Results}
\begin{itemize}
  \item object detection baselines: RALA, Ours, RMT-T.
  \item image classification baselines: RALA, Ours, MLLA-T.
  \item inference speed
\end{itemize}

% \subsection{Datasets}

% \subsection{Evaluation Protocol}
% This subsection defines how model performance is assessed. It should include:
% \begin{itemize}
%   \item Evaluation metrics (\eg, accuracy, AUROC, F1-score, mAP).
%   \item Data splitting strategy and any cross-validation or repeated-trial setup.
%   \item Baselines used for comparison and criteria for fair evaluation.
%   \item Any task-specific rules or conventions followed during evaluation.
% \end{itemize}

% \subsection{Experimental Settings}
% This subsection outlines the full implementation and training configuration for reproducibility:
% \begin{itemize}
%   \item Model architecture settings and important hyperparameters
%         (\eg, learning rate, batch size, number of epochs).
%   \item Optimization algorithms and loss functions.
%   \item Hardware and software environment (\eg, GPU type, deep learning framework).
%   \item Data preprocessing and augmentation pipeline.
% \end{itemize}

% \subsection{Results and Discussion}
% This subsection presents and analyzes the experimental results:
% \begin{itemize}
%   \item Quantitative results summarized in tables and figures.
%   \item Qualitative results when appropriate.
%   \item Comparisons with state-of-the-art methods.
%   \item Insights on why the proposed method performs well and any limitations or failure cases.
% \end{itemize}

% \subsubsection{Comparison Experiments}
% This subsection reports direct comparisons with existing approaches:
% \begin{itemize}
%   \item Performance comparison with classical baselines and recent state-of-the-art methods.
%   \item Analysis of cases where the proposed method excels.
%   \item Discussion of scenarios where performance is comparable or lower.
%   \item Positioning of the proposed method within the broader research landscape.
%         There should be multiple tables \ref{tab1} and figures \ref{fig2} for demonstration and clarity.
% \end{itemize}

% \begin{table}[h]
%   \centering
%   \caption{Experimental results of predicate classification without Graph Constraint.}
%   \label{table:pred_cls_results}
%   \begin{tabular}{ |c|c c c| }
%     \hline
%     Category & R@20   & R@50   & R@100  \\
%     \hline
%     Spatial  & 0.4471 & 0.6356 & 0.7590 \\
%     % \hline
%     Action   & 0.2971 & 0.3835 & 0.4544 \\
%     % \hline
%     $mR@K$   & 0.0596 & 0.1027 & 0.1530 \\
%     % \hline
%     $R@K$    & 0.2991 & 0.4177 & 0.4987 \\
%     \hline
%   \end{tabular}
% \end{table}


% \begin{table}[t]
%   \centering
%   \caption{Table captions should be placed above the tables. Tables should be showned on top. Best values are bolded while second-best ones are underlined.}\label{tab1}
%   \begin{tabular}{|p{3cm}|c|c|}
%     \hline
%     Methods       & Metric 1          & Metric 2          \\
%     \hline
%     SOTA 1 [ref]  & 0.001             & 0.002             \\
%     SOTA 2 [ref]  & 0.003             & \underline{0.008} \\
%     SOTA 3 [ref]  & 0.005             & 0.006             \\
%     SOTA 4 [ref]  & \underline{0.007} & 0.004             \\
%     \hline
%     \textbf{Ours} & \textbf{0.009}    & \textbf{0.010}    \\
%     \hline
%   \end{tabular}
% \end{table}

% \begin{figure}[t]
%   \includegraphics[width=\textwidth]{figures/fig1.eps}
%   \caption{A figure caption is always placed below the illustration.
%     Please note that short captions are centered, while long ones are
%     justified by the macro package automatically.} \label{fig2}
% \end{figure}

% \subsubsection{Ablation Study}
% This subsection evaluates the contribution of each component of the proposed method:
% \begin{itemize}
%   \item Removal or modification of individual modules or strategies.
%   \item Measurement of performance change to quantify each component’s importance.
%   \item Justification of design decisions based on empirical evidence.
% \end{itemize}

\begin{table}[htbp]
  \centering
  \begin{tabular}{c|c|cc|c}
    \hline
    Cost & Model & Params(M) & FLOPs(G) & Top1-acc(\%) \\
    \hline
    \multirow{3}{*}{\rotatebox[origin=c]{90}{\makecell{$\sim 2.5$G}}}
         & RALA  &           &          & \textbf{}    \\
         & RMT-T &           &          & \textbf{}    \\
         & Ours  &           &          & \textbf{}    \\
    \hline
  \end{tabular}
  \vspace{2mm}
  \caption{Classification results of different general backbones on ImageNet-100}
  \label{tab:classification}
\end{table}

\begin{table}[htbp]
  \centering
  \begin{tabular}{c|c|cc|c}
    \hline
    Cost & Model  & Params(M) & FLOPs(G) & Top1-acc(\%) \\
    \hline
    \multirow{3}{*}{\rotatebox[origin=c]{90}{\makecell{$\sim 2.5$G}}}
         & RALA   &           &          & \textbf{}    \\
         & MLLA-T &           &          & \textbf{}    \\
         & Ours   &           &          & \textbf{}    \\
    \hline
  \end{tabular}
  \vspace{2mm}
  \caption{Object detection results of different general backbones on COCO 2017}
  \label{tab:object_detection}
\end{table}

\subsection{Ablation Study}
\subsubsection{Other gating mechanism}
Discuss the result of other gating mechanism compare to RL.
We hypothesize that the RL gating will be superior to the existing methods due to 
its ability to gather global information. On the other hand, the 
existing methods may have better speed, but it should be negligible.
\subsubsection{Without budget concern}
Discuss the result fo RL gating without budget concern. We hypothesize that the RL will be
lazy to assign a low weight for tokens, since it only needs to retain the performance. 
\subsubsection{Without Fisher Initialization}
Discuss the training time, result without Fisher Initialzatio for Actor-Critic network. We propose 
that training the Actor-Critic network will be difficult due to unstablized training, and the high variance
of RL. 
% explain why it works or not?
% what's next?
